% PAGES: 257-258

Also, let $\widehat\sigma_{Y,\bfx}$ be the sample covariance vector of $Y$ and $X_1,
\ldots, X_n$ given by
\begin{equation}
\widehat\sigma_{Y,\bfx} \;=\; \begin{pmatrix} \widehat{\cov}(Y,X_1) \\ \vdots \\
\widehat{\cov}(Y,X_n) \end{pmatrix}.
\end{equation}

Let $\overline{T}_Y$ and $\overline{T}_{X_k}$ be the column vectors of the time series
data for the random variables $Y$ and $X_k$ normalized to mean equal to $0$, for
$k=1:n$, respectively, i.e.,
\begin{align}
\overline{T}_Y &= (Y(t_i)-\widehat\mu_Y)_{i=1:N} \;=\; T_Y - \widehat\mu_Y\bfone;\\
\overline{T}_{X_k} &= (X_k(t_i)-\widehat\mu_{X_k})_{i=1:N} \;=\; T_{X_k} -
\widehat\mu_{X_k}\bfone,
\end{align}
and let
\begin{equation}
\overline{T}_{\bfx} \;=\; col\left(\overline{T}_{X_k}\right)_{k=1:n}.
\end{equation}

\begin{lemma}
Let $Y$ and $X_1, X_2, \ldots, X_n$ be random variables given by the $N\times1$ time
series data vectors $T_Y$ and $T_{X_1}, T_{X_2}, \ldots, T_{X_n}$, respectively. Let
$T_{\bfx} = col(T_{X_k})_{k=1:n}$. Assume that the sample covariance matrix
$\widehat\Sigma_{\bfx}$ corresponding to the time series data for $X_1, X_2, \ldots,
X_n$ is nonsingular.

The solution $a$ and $\bfb = (b_i)_{i=1:n}$ to the linear regression problem
\begin{align}
T_Y &\approx (\bfone\ \ T_{X_1}\ \ \ldots\ \ T_{X_n})\begin{pmatrix} a \\ \bfb
\end{pmatrix}\\
&= a\bfone + b_1T_{X_1} + b_2T_{X_2} + \ldots + b_nT_{X_n}, \notag
\end{align}
where $\bfone$ is the $N\times1$ column vector with all entries equal to $1$, is
\begin{align}
\bfb &= \left(\overline{T}_{\bfx}^t\overline{T}_{\bfx}\right)^{-1}\overline{T}_{\bfx}^t
\overline{T}_Y;\\
a &= \widehat\mu_Y - \widehat\mu_{\bfx}^t\left(\overline{T}_{\bfx}^t
\overline{T}_{\bfx}\right)^{-1}\overline{T}_{\bfx}^t\overline{T}_Y,
\end{align}
where $\overline{T}_Y$ and $\overline{T}_{\bfx}$ are given by (8.74) and (8.76),
respectively.

Moreover, (8.78) and (8.79) can be written as
\begin{align}
\bfb &= \left(\widehat\Sigma_{\bfx}\right)^{-1}\widehat\sigma_{Y,\bfx}.\\
a &= \widehat\mu_Y - \widehat\mu_{\bfx}^t\left(\widehat\Sigma_{\bfx}\right)^{-1}
\widehat\sigma_{Y,\bfx},
\end{align}
where $\widehat\mu_Y$, $\widehat\mu_{\bfx}$, and $\widehat\sigma_{Y,\bfx}$ are given by
(8.70), (8.72), and (8.73), respectively, and $\widehat\Sigma_{\bfx}$ is the sample
covariance matrix given by (7.42).
\end{lemma}

The results of Lemma 8.2 connect the linear regression ordinary least squares solution
for time series data to the ordinary least squares solution for random variables from
section 8.3; see (8.57) and (8.58).

\begin{proof}
% The "Since ..." clause below is printed with an overline, $\overline{T}_\bfx$, even
% though the equation (8.82) it introduces uses the plain (non-overlined) $T_\bfx =
% col(T_{X_k})_{k=1:n}$ notation for the ORIGINAL, non-mean-centered data -- matching
% (8.77) and the $T_\bfx$ notation from section 8.4's opening. This mismatch is exactly
% as printed on the page (verified at 1200 dpi); transcribed as-is per the fidelity
% rule rather than silently "corrected."
Since $\overline{T}_{\bfx} = col\left(\overline{T}_{X_k}\right)_{k=1:n}$, the ordinary
least squares problem (8.77) can also be written as follows:
\begin{equation}
T_Y \;\approx\; (\bfone\ \ T_{\bfx})\begin{pmatrix} a \\ \bfb \end{pmatrix}.
\end{equation}

Recall from section 8.1 that the solution to the least squares problem $y\approx Ax$
requires solving the linear system $A^ty = A^tAx$. Problem (8.82) corresponds to
\[
y \approx Ax \quad \text{with} \quad y = T_Y,\ A = (\bfone\ \ T_{\bfx}),\ x =
\begin{pmatrix} a \\ \bfb \end{pmatrix},
\]
and therefore $A^ty = A^tAx$ can be written as
\begin{equation}
\begin{pmatrix} \bfone^t \\ T_{\bfx}^t \end{pmatrix}T_Y \;=\; \begin{pmatrix} \bfone^t \\
T_{\bfx}^t \end{pmatrix}(\bfone\ \ T_{\bfx})\begin{pmatrix} a \\ \bfb \end{pmatrix}.
\end{equation}

By looking only at the first row of (8.83), we obtain that
\begin{align}
\bfone^tT_Y &= \bfone^t(\bfone\ \ T_{\bfx})\begin{pmatrix} a \\ \bfb \end{pmatrix} \;=\;
\bfone^t(\bfone a + T_{\bfx}\bfb) \notag\\
&= \bfone^t\bfone a + \bfone^tT_{\bfx}\bfb.
\end{align}

From (8.70), we find that
\begin{equation}
\bfone^tT_Y \;=\; N\widehat\mu_Y.
\end{equation}

Also, note that
\begin{equation}
\bfone^t\bfone \;=\; \sum_{i=1}^N 1 \;=\; N.
\end{equation}

Recall that $T_{\bfx} = col(T_{X_k})_{k=1:n}$. From (8.71), it follows that
$\bfone^tT_{X_k} = N\widehat\mu_{X_k}$ for all $k=1:n$, and therefore $\bfone^tT_{\bfx}$
is the following row vector:
\begin{align}
\bfone^tT_{\bfx} &= \left(\bfone^tT_{X_1}\ \ \bfone^tT_{X_2}\ \ \ldots\ \
\bfone^tT_{X_n}\right) \notag\\
&= \left(N\widehat\mu_{X_1}\ \ N\widehat\mu_{X_2}\ \ \ldots\ \ N\widehat\mu_{X_n}\right)
\notag\\
&= N\left(\widehat\mu_{X_1}\ \ \widehat\mu_{X_2}\ \ \ldots\ \ \widehat\mu_{X_n}\right)
\notag\\
&= N\widehat\mu_{\bfx}^t,
\end{align}
where $\widehat\mu_{\bfx}^t = \left(\widehat\mu_{X_1}\ \ \widehat\mu_{X_2}\ \ \ldots\ \
\widehat\mu_{X_n}\right)$; see (8.72).

From (8.85), (8.86), and (8.87), we obtain that (8.84) can be written as
\begin{equation}
N\widehat\mu_Y \;=\; Na + N\widehat\mu_{\bfx}^t\bfb.
\end{equation}

By solving (8.88) for $a$, we obtain that
\begin{equation}
a \;=\; \widehat\mu_Y - \widehat\mu_{\bfx}^t\bfb.
\end{equation}

Note that the ordinary least squares problem (8.82) can be written as
\begin{equation}
T_Y \;\approx\; (\bfone\ \ T_{\bfx})\begin{pmatrix} a \\ \bfb \end{pmatrix} \;=\;
a\bfone + T_{\bfx}\bfb.
\end{equation}

By substituting the expression (8.89) for $a$ into (8.90), we find that
\begin{align}
T_Y &\approx a\bfone + T_{\bfx}\bfb \notag\\
&\iff T_Y \approx \widehat\mu_Y\bfone - \left(\widehat\mu_{\bfx}^t\bfb\right)\bfone +
T_{\bfx}\bfb \notag\\
&\iff T_Y - \widehat\mu_Y\bfone \approx T_{\bfx}\bfb -
\left(\widehat\mu_{\bfx}^t\bfb\right)\bfone.
\end{align}
